\documentclass[conference]{IEEEtran}
\IEEEoverridecommandlockouts
\usepackage[utf8]{inputenc}   % Türkçe karakterler
\usepackage[T1]{fontenc}      % ç ş ğ ı ö ü İ
\usepackage{cite}
\usepackage{amsmath,amssymb,amsfonts}
\usepackage{algorithmic}
\usepackage{graphicx}
\usepackage{textcomp}
\usepackage{xcolor}
\usepackage{booktabs}
\usepackage{url}
\begin{document}

\title{Film Afişlerinden Çoklu-Etiketli Tür Tahmini:\\
Beş Görsel Transformatör Modelinin Karşılaştırmalı Değerlendirmesi}

\author{\IEEEauthorblockN{Mustafa Kerem Çekici}
\IEEEauthorblockA{\textit{Bilişim Sistemleri Mühendisliği Bölümü} \\
\textit{Kocaeli Üniversitesi}\\
İzmit, Türkiye \\
231307121 -- keremcek70@gmail.com}
}

\maketitle

\begin{abstract}
Bu çalışmada, film afişi görselinden filmin türlerinin tahmin edilmesi
çoklu-etiketli (multi-label) bir sınıflandırma problemi olarak ele alınmıştır.
TMDB'nin resmî API'si kullanılarak 15 hedef tür için eksiklik-güdümlü bir
örnekleme algoritmasıyla dengeli ve kombinasyon-farkında bir veri kümesi
(23.640 film) oluşturulmuştur. Beş farklı görsel transformatör modeli
(ViT, DeiT, BeiT, Swin ve CvT) ImageNet ön-eğitiminden başlanarak ince
ayarlanmış ve 5-katlı çapraz doğrulama ile değerlendirilmiştir. En iyi sonucu
Swin modeli vermiştir (makro F1 = 0,562; makro AUC = 0,862), bu da ön-eğitimli
CNN tabanlı bir referansa (0,383) göre +0,179 ve şans seviyesindeki rasgele bir
tahmin ediciye (0,149) göre yaklaşık 3,8 katlık bir iyileşmedir. Dengeli veri
sayesinde önceki çalışmadaki aşırı tahmin sorunu giderilmiş (film başına ortalama
$\sim$2,4 tür, gerçek değer 2,18) ve frekans biası ortadan kalkmıştır; nadir
türlerdeki belirgin başarı artışı modelin görsel sinyali gerçekten öğrendiğini,
salt etiket korelasyonlarını ezberlemediğini göstermektedir.
\end{abstract}

\begin{IEEEkeywords}
çoklu-etiket sınıflandırma, görsel transformatör, film afişi, transfer öğrenme,
çapraz doğrulama, TMDB
\end{IEEEkeywords}

\section{Giriş}
Film afişleri, türü tek bir bakışta sezdiren güçlü görsel kodlar içerir: korku
afişlerindeki karanlık tonlar, animasyonun çizgi-film estetiği, romantik
filmlerin sıcak renk paleti gibi. Bu çalışmanın amacı, yalnızca afiş görselini
girdi alarak bir filmin türlerini tahmin eden bir model geliştirmektir. Bir film
birden çok türe ait olabildiğinden problem doğal olarak \textit{çoklu-etiketlidir}.

İlk denememizde (v1) sıfırdan eğitilen bir evrişimli ağ (CNN) ve dondurulmuş bir
MobileNetV3 referansı kullanılmış; ancak dengesiz veri yüzünden model türlerin
görsel özelliklerini öğrenmek yerine sınıf frekanslarını ezberlemiş, film başına
5--7 tür tahmin ederek aşırı tahmin yapmıştır. Bu rapor, problemi sıfırdan ele
alan ikinci sürümü (v2) sunar. Katkılarımız:
\begin{itemize}
\item TMDB API üzerinden \textbf{dengeli ve kombinasyon-farkında} bir veri kümesi
oluşturan iki aşamalı bir örnekleme algoritması;
\item Beş farklı görsel transformatörün ($ViT$, $DeiT$, $BeiT$, $Swin$, $CvT$)
aynı koşullarda 5-katlı çapraz doğrulamayla karşılaştırılması;
\item Aşırı tahmin ve frekans biası sorunlarının nicel olarak giderildiğinin
gösterilmesi ve rasgele/frekans-öncül referanslarla modelin gerçek görsel öğrenme
yaptığının kanıtlanması.
\end{itemize}

\section{Veri Toplama ve Temizleme}
\subsection{HTML Kazımadan Resmî API'ye Geçiş}
İlk yaklaşımda film verileri \texttt{themoviedb.org} sitesinin HTML sayfaları
kazınarak toplanmıştır. Site agresif hız sınırlaması uyguladığından, TLS parmak
izi taklidi ve VPN kullanılmasına rağmen sürekli HTTP 429 (Too Many Requests)
yanıtları alınmıştır: tek bir hedefli kazıma oturumunda 9.619 adet 429 yanıtı
kaydedilmiş ve 864 film kalıcı olarak kaybedilmiştir; pratik hız 3 saatte
$\sim$500 film düzeyinde kalmıştır.

Bu nedenle veri toplama, TMDB'nin \textbf{resmî JSON API}'sine taşınmıştır. API
kâr amacı gütmeyen kullanımda ücretsizdir, günlük istek limiti yoktur ve
$\sim$40--50 istek/saniyeye izin verir. \texttt{/discover/movie} uç noktası, her
filmin tür listesini (\texttt{genre\_ids}) ve afiş yolunu (\texttt{poster\_path})
tek istekte döndürdüğünden, çoklu-etiketler ek maliyet olmadan elde edilmiş ve
afişler yalnızca kabul edilen filmler için indirilmiştir.

\subsection{Dengeli ve Kombinasyon-Farkında Örnekleme}
v1 veri kümesi ileri derecede dengesizdi (Drama 6.381 örnek, History 858 örnek;
$\sim$7,4 kat). Dengeyi sağlamak için iki aşamalı bir yöntem kullanılmıştır:
\textbf{(1) Aday havuzu:} her tür için \texttt{with\_genres} filtresiyle
\texttt{vote\_count}$\geq$10 olan aday filmler toplanır (kalitesiz/amatör/postersiz
girdiler elenir). \textbf{(2) Eksiklik-güdümlü açgözlü seçim:} her adımda en eksik
tür belirlenir ve bu türü içeren, en az sayıda tür taşıyan film seçilir; baskın
türler (Drama, Comedy) ``yolcu'' olarak en sona bırakılır. Böylece nadir türler
hedefe ulaşırken baskın türlerin aşırı şişmesi sınırlanır.

\texttt{vote\_count} eşiği bir kalite ölçüsü değil, bilinirlik göstergesidir;
ancak filtresiz çekim çöp doludur (örn. Documentary için $\sim$220.000 girdi).
Eşik 30 alındığında History yalnızca 2.149 filmde kalıyordu; eşik 10'a
düşürüldüğünde tüm nadir türler 3.000 hedefine ulaşabilecek kadar
($\geq$3.850) aday sağlamıştır.

\subsection{Sonuç Veri Kümesi}
Nihai veri kümesi \textbf{23.640 film} içerir; film başına ortalama 2,18 tür
düşmektedir ve afiş kapsamı \%100'dür. Dengesizlik 7,4 kattan 2,22 kata
düşürülmüştür (nadir/orta türlerin tamamı $\sim$3.000 örneğe çıkarılmış,
History 858'den 3.000'e taşınmıştır). Çekim öncesi/sonrası tür dağılımı
Şekil~\ref{fig:dist}'te, türler arası birlikte-geçiş (co-occurrence) matrisi
Şekil~\ref{fig:cooc}'da verilmiştir.

\begin{figure}[htbp]
\centerline{\includegraphics[width=\columnwidth]{dist_before_after.png}}
\caption{Tür dağılımı: çekim öncesi (v1, dengesiz) ve sonrası (v2, dengeli).}
\label{fig:dist}
\end{figure}

\begin{figure}[htbp]
\centerline{\includegraphics[width=\columnwidth]{cooccurrence_v2.png}}
\caption{v2 veri kümesinde türler arası birlikte-geçiş matrisi.}
\label{fig:cooc}
\end{figure}

\section{Yöntem}
\subsection{Modeller}
Beş farklı görsel transformatör, ImageNet ön-eğitimli ağırlıklardan başlanarak
ince ayarlanmıştır: \textbf{ViT} (saf öz-dikkat tabanlı transformatör)~\cite{vit},
\textbf{DeiT} (veri-verimli, damıtma ile eğitilen ViT)~\cite{deit}, \textbf{BeiT}
(maskeli görüntü modellemesiyle ön-eğitilmiş)~\cite{beit}, \textbf{Swin}
(kaydırmalı pencerelerle hiyerarşik transformatör)~\cite{swin} ve \textbf{CvT}
(evrişim katmanlarıyla zenginleştirilmiş transformatör)~\cite{cvt}. Tüm modeller
\textit{base} ölçeğinde olup CvT-13 daha küçük kapasitelidir. Modeller HuggingFace
Transformers~\cite{hf} kütüphanesi üzerinden, sınıflandırıcı başı 15 türe
uyarlanarak (\texttt{problem\_type = multi\_label\_classification}) yüklenmiştir.

\subsection{Ön İşleme ve Eğitim}
Afişler 2:3 dikey orana sahiptir; her model için $224\times224$ piksele yeniden
boyutlandırılmış (kareye sıkıştırma), modelin kendi ortalama/standart sapma
değerleriyle normalize edilmiş ve eğitimde hafif renk artırımı uygulanmıştır.
Kayıp fonksiyonu \texttt{BCEWithLogitsLoss}'tur; veri dengeli olduğundan ağır
\texttt{pos\_weight} kullanılmamıştır (v1'de bu, aşırı tahmine yol açmıştı).
Optimizasyon AdamW (öğrenme oranı $5\times10^{-5}$, ağırlık sönümü $0{,}01$),
kosinüs ısınma çizelgesi ve karma duyarlık (AMP) ile yapılmıştır; toplu boyut 64,
epoch sayısı 5'tir.

\subsection{Çapraz Doğrulama ve Değerlendirme}
Veri kümesi, çoklu-etiket dağılımını koruyan \texttt{MultilabelStratifiedKFold}
ile~\cite{stratif} \textbf{5 kata} bölünmüştür (kat başına $\sim$18.900 eğitim /
$\sim$4.730 doğrulama; katlar arası örtüşme sıfır, mükemmel katmanlaşma). Her
örnek tam olarak bir kez doğrulama katında bulunduğundan, beş katın doğrulama
tahminleri birleştirilerek \textbf{kat-dışı (out-of-fold, OOF)} tahminler elde
edilmiş ve tüm metrikler bu birleşik tahminler üzerinden hesaplanmıştır. Sınıf
başına eşik, F1'i en üst düzeye çıkaracak biçimde $[0{,}3;\,0{,}7]$ aralığında
optimize edilmiştir. Çoklu-etiket olduğundan Accuracy, Precision,
Recall (Sensitivity), Specificity, F-Score ve AUC değerleri sınıf başına
hesaplanıp makro/mikro ortalaması alınmıştır.

\section{Deneysel Sonuçlar}
\subsection{Model Karşılaştırması}
Tablo~\ref{tab:models} beş modelin tam metrik setini v1 referansları ve rasgele
referanslarla birlikte sunar. \textbf{En iyi model Swin}'dir (makro F1 = 0,562).
Beş transformatörün tamamı v1 referansını (0,383) açık biçimde geçmiştir.

\begin{table*}[htbp]
\caption{Model karşılaştırması (OOF, eşik-optimize)}
\begin{center}
\small
\begin{tabular}{lcccccccc}
\toprule
\textbf{Model} & \textbf{Makro F1} & \textbf{Mikro F1} & \textbf{Makro AUC} & \textbf{Prec.} & \textbf{Rec.} & \textbf{Makro Acc.} & \textbf{Makro Spec.}\\
\midrule
\textbf{Swin} & \textbf{0,562} & 0,566 & 0,862 & 0,541 & 0,587 & \textbf{0,867} & \textbf{0,912}\\
BeiT & 0,534 & 0,539 & 0,844 & 0,517 & 0,555 & 0,860 & 0,908\\
ViT  & 0,527 & 0,535 & 0,841 & 0,512 & 0,545 & 0,860 & 0,909\\
DeiT & 0,524 & 0,530 & 0,836 & 0,519 & 0,532 & 0,860 & 0,912\\
CvT  & 0,501 & 0,509 & 0,823 & 0,474 & 0,536 & 0,846 & 0,895\\
\midrule
Frozen baseline (v1) & 0,383 & -- & -- & -- & -- & -- & --\\
Scratch CNN (v1) & 0,333 & -- & -- & -- & -- & -- & --\\
Rasgele ($\sim$2,4/film) & 0,149 & -- & -- & -- & -- & -- & --\\
Frekans-öncül & 0,146 & -- & -- & -- & -- & -- & --\\
\bottomrule
\end{tabular}
\label{tab:models}
\end{center}
\end{table*}

Sınıf başına F1 değerlerinin model bazında karşılaştırması
Şekil~\ref{fig:perclass}'te verilmiştir.

\begin{figure*}[htbp]
\centerline{\includegraphics[width=\textwidth]{per_class_f1_models.png}}
\caption{Sınıf başına F1 değerleri (beş model).}
\label{fig:perclass}
\end{figure*}

\subsection{En İyi Modelin Sınıf Bazlı Metrikleri}
Tablo~\ref{tab:perclass} Swin modelinin sınıf başına tüm metriklerini gösterir.
En yüksek başarı Animation türündedir (F1 = 0,863; AUC = 0,980) -- çizgi-film
estetiği güçlü ve ayırt edici bir görsel imza taşır. En düşük başarılar
görsel olarak en belirsiz türlerde görülmektedir (Fantasy 0,419; Mystery 0,432;
Crime 0,440). Özgüllük (specificity) değerlerinin yüksekliği (0,78--0,99) yanlış
pozitiflerin az olduğunu, yani aşırı tahminin olmadığını doğrular.

\begin{table*}[htbp]
\caption{Swin modeli -- sınıf başına metrikler (OOF, eşik-optimize)}
\begin{center}
\begin{tabular}{lcccccc}
\toprule
\textbf{Tür} & \textbf{Accuracy} & \textbf{Precision} & \textbf{Recall/Sens.} & \textbf{Specificity} & \textbf{F1} & \textbf{AUC}\\
\midrule
Action          & 0,874 & 0,596 & 0,643 & 0,917 & 0,619 & 0,890\\
Adventure       & 0,850 & 0,422 & 0,478 & 0,905 & 0,448 & 0,821\\
Animation       & 0,966 & 0,896 & 0,833 & 0,986 & \textbf{0,863} & 0,980\\
Comedy          & 0,869 & 0,620 & 0,669 & 0,912 & 0,644 & 0,887\\
Crime           & 0,847 & 0,411 & 0,473 & 0,901 & 0,440 & 0,816\\
Documentary     & 0,893 & 0,573 & 0,614 & 0,933 & 0,593 & 0,901\\
Drama           & 0,753 & 0,550 & 0,674 & 0,784 & 0,606 & 0,811\\
Family          & 0,905 & 0,613 & 0,681 & 0,938 & 0,645 & 0,917\\
Fantasy         & 0,850 & 0,413 & 0,426 & 0,912 & 0,419 & 0,782\\
History         & 0,863 & 0,461 & 0,482 & 0,918 & 0,471 & 0,841\\
Horror          & 0,901 & 0,601 & 0,662 & 0,936 & 0,630 & 0,910\\
Mystery         & 0,848 & 0,412 & 0,454 & 0,906 & 0,432 & 0,803\\
Romance         & 0,875 & 0,506 & 0,563 & 0,920 & 0,533 & 0,866\\
Science Fiction & 0,889 & 0,564 & 0,559 & 0,937 & 0,561 & 0,868\\
Thriller        & 0,827 & 0,475 & 0,590 & 0,873 & 0,526 & 0,844\\
\midrule
\textbf{Makro}  & 0,867 & 0,541 & 0,587 & 0,912 & \textbf{0,562} & 0,862\\
\bottomrule
\end{tabular}
\label{tab:perclass}
\end{center}
\end{table*}

\subsection{Karmaşıklık Matrisleri ve ROC Eğrileri}
Beş modelin tamamı için sınıf başına karmaşıklık matrisleri ve teke-karşı-tümü
(one-vs-rest) ROC eğrileri üretilmiştir; yer kısıtı nedeniyle en iyi modelin
(Swin) sonuçları Şekil~\ref{fig:cm} ve Şekil~\ref{fig:roc}'da gösterilmektedir.
Diğer dört modelin (ViT, DeiT, BeiT, CvT) karmaşıklık matrisleri ve ROC eğrileri
üretilmiş olup yer kısıtı nedeniyle rapora eklenmemiştir.

\begin{figure}[htbp]
\centerline{\includegraphics[width=\columnwidth]{confusion_swin.png}}
\caption{Swin -- sınıf başına karmaşıklık matrisleri.}
\label{fig:cm}
\end{figure}

\begin{figure}[htbp]
\centerline{\includegraphics[width=\columnwidth]{roc_swin.png}}
\caption{Swin -- ROC eğrileri (one-vs-rest), makro AUC = 0,862.}
\label{fig:roc}
\end{figure}

\subsection{Eğitim Eğrileri ve Süreler}
Her modelin beş katına ait eğitim/doğrulama kayıp eğrileri
Şekil~\ref{fig:loss}'te verilmiştir. Yüksek kapasiteli modeller (ViT, DeiT,
BeiT, Swin) 2.--3. epoch'tan sonra aşırı öğrenmeye başlarken (eğitim kaybı
düşerken doğrulama kaybı yükselir), CvT 5 epoch'ta hâlâ yetersiz öğrenme
(underfitting) bölgesindedir; bu, kapasite/veri dengesinin bir göstergesidir.
En iyi-F1 kontrol noktası saklandığından aşırı öğrenme son modeli olumsuz
etkilememiştir. Eğitim ve çıkarım süreleri Tablo~\ref{tab:time}'de özetlenmiştir.

\begin{figure*}[htbp]
\centerline{\includegraphics[width=\textwidth]{loss_curves.png}}
\caption{Kat başına eğitim/doğrulama kayıp eğrileri (her model).}
\label{fig:loss}
\end{figure*}

\begin{table}[htbp]
\caption{Eğitim/çıkarım süresi ve ortalama tahmin sayısı}
\begin{center}
\begin{tabular}{lccc}
\toprule
\textbf{Model} & \textbf{Eğitim (dk/kat)} & \textbf{Çıkarım (ms/görsel)} & \textbf{Ort. tür/film}\\
\midrule
ViT  & 3,94 & 1,58 & 2,36\\
DeiT & 3,96 & 1,58 & 2,29\\
BeiT & 4,16 & 1,60 & 2,39\\
Swin & 5,84 & 1,63 & 2,41\\
CvT  & 4,73 & 1,58 & 2,52\\
\bottomrule
\end{tabular}
\label{tab:time}
\end{center}
\end{table}

\subsection{Kritik Kontroller}
\textbf{Aşırı tahmin giderildi:} tüm modeller film başına ortalama 2,3--2,5 tür
tahmin etmektedir (gerçek değer 2,18), oysa v1'de bu 5--7 idi.
\textbf{Frekans biası giderildi:} v1'de en başarılı türler en sık görülen
Drama/Comedy, en başarısızları ise nadir History/Documentary ($\sim$0,10) idi.
v2'de en başarılı tür Animation (0,863) olmuş; History 0,10'dan 0,47'ye,
Documentary 0,11'den 0,59'a yükselmiştir. Bu, modelin frekans yerine görsel
sinyali öğrendiğini gösterir.

\textbf{Rasgele referansla karşılaştırma:} film başına $\sim$2,4 tür rasgele
tahmin edildiğinde makro F1 yalnızca 0,149; afişe hiç bakmadan sınıf
frekanslarına göre tahmin (frekans-öncül) edildiğinde 0,146 elde edilmektedir.
Swin'in 0,562 değeri, şans seviyesinin \textbf{$\sim$3,8 katıdır} ve
frekans-öncül referansı 0,42 puan geçmektedir; dolayısıyla tahmin gücü
korelasyon/frekanstan değil, doğrudan afiş görselinden gelmektedir. Örnek
tahminler Şekil~\ref{fig:samples}'te sunulmuştur.

\begin{figure}[htbp]
\centerline{\includegraphics[width=\columnwidth]{samples_swin.png}}
\caption{Swin -- örnek afişler için gerçek ve tahmin edilen türler.}
\label{fig:samples}
\end{figure}

\section{Tartışma}
Sonuçlar, dengeli ve kombinasyon-farkında veri toplamanın çoklu-etiket
afiş sınıflandırmasında belirleyici olduğunu göstermektedir: v1'deki aşırı
tahmin ve frekans biası, ağır kayıp ağırlıklandırmasından değil, veri
dengesizliğinden kaynaklanıyordu ve kaynağında çözülmüştür. Mimari açıdan,
hiyerarşik ve yerel-dikkatli Swin'in en iyi sonucu vermesi, görece küçük
veri kümelerinde tümevarımsal ön-yargısı daha güçlü modellerin saf ViT'e
üstün geldiği yönündeki bulgularla tutarlıdır. CvT'nin düşük başarısı, küçük
kapasitesi ve 5 epoch'ta yetersiz öğrenmesiyle açıklanabilir.

\textbf{Co-occurrence (etiket korelasyonu) riski:} Comedy--Romance gibi sık
birlikte görülen türlerde modelin görsel öğrenme yerine etiket korelasyonunu
sömürme riski vardır. Veri kümesini kombinasyon dengesine dikkat ederek
oluşturmak bu riski azaltmıştır ve nadir türlerdeki büyük başarı artışı riskin
baskın olmadığını göstermektedir; ancak bu risk tamamen elenmiş değildir.
Ayrıca afişlerin kareye sıkıştırılması en/boy oranını bozmaktadır; en zayıf
türler (Fantasy, Mystery, Crime) görsel olarak en belirsiz türlerdir.

\section{Sonuç}
Film afişinden çoklu-etiketli tür tahmini için beş görsel transformatör aynı
koşullarda karşılaştırılmıştır. En iyi model Swin, makro F1 = 0,562 ile v1
referansına göre yaklaşık \%47 görece iyileşme sağlamış; aşırı tahmin ve frekans
biası nicel olarak giderilmiş, rasgele referanslarla modelin görsel sinyali
gerçekten öğrendiği kanıtlanmıştır. Gelecek çalışmalarda en/boy oranını koruyan
ön işleme, daha güçlü düzenlileştirme ve metin (başlık) ile çok-modlu birleşim
incelenebilir.

\vspace{4pt}
\noindent\footnotesize{%
$^{1}$\textit{Kod ve notebook'lar (GitHub):}
\url{https://github.com/Sayicon/poster-based-genre-prediction}.
\textit{Eğitim (09):}
\url{https://colab.research.google.com/github/Sayicon/poster-based-genre-prediction/blob/main/notebooks/09_train_transformer.ipynb}.
\textit{Değerlendirme (10):}
\url{https://colab.research.google.com/github/Sayicon/poster-based-genre-prediction/blob/main/notebooks/10_evaluate.ipynb}.%
}

\begin{thebibliography}{00}
\bibitem{vit} A. Dosovitskiy \textit{et al.}, ``An image is worth 16x16 words: Transformers for image recognition at scale,'' in \textit{Proc. ICLR}, 2021.
\bibitem{deit} H. Touvron \textit{et al.}, ``Training data-efficient image transformers \& distillation through attention,'' in \textit{Proc. ICML}, 2021.
\bibitem{beit} H. Bao, L. Dong, S. Piao, and F. Wei, ``BEiT: BERT pre-training of image transformers,'' in \textit{Proc. ICLR}, 2022.
\bibitem{swin} Z. Liu \textit{et al.}, ``Swin transformer: Hierarchical vision transformer using shifted windows,'' in \textit{Proc. ICCV}, 2021, pp. 10012--10022.
\bibitem{cvt} H. Wu \textit{et al.}, ``CvT: Introducing convolutions to vision transformers,'' in \textit{Proc. ICCV}, 2021, pp. 22--31.
\bibitem{stratif} K. Sechidis, G. Tsoumakas, and I. Vlahavas, ``On the stratification of multi-label data,'' in \textit{Proc. ECML PKDD}, 2011, pp. 145--158.
\bibitem{hf} T. Wolf \textit{et al.}, ``Transformers: State-of-the-art natural language processing,'' in \textit{Proc. EMNLP: System Demonstrations}, 2020, pp. 38--45.
\bibitem{tmdb} The Movie Database (TMDB), ``API documentation,'' \url{https://developer.themoviedb.org}, erişim: 2026.
\end{thebibliography}

\end{document}
